\documentclass[lettersize,journal]{IEEEtran}

\usepackage{amsmath,amsfonts}
\usepackage{array}
\usepackage[caption=false,font=normalsize,labelfont=sf,textfont=sf]{subfig}
\usepackage{textcomp}
\usepackage{stfloats}
\usepackage{url}
\usepackage{graphicx}
\usepackage{cite}
\usepackage{xcolor}
\hyphenation{op-tical net-works semi-conduc-tor IEEE-Xplore}
\def\BibTeX{{\rm B\kern-.05em{\sc i\kern-.025em b}\kern-.08em
		T\kern-.1667em\lower.7ex\hbox{E}\kern-.125emX}}
\usepackage{balance}
\usepackage{booktabs}
\usepackage{tabularx}
\usepackage{ragged2e}
\usepackage{float}
\usepackage{algorithm}
\usepackage{algpseudocode}


\hyphenation{op-tical net-works semi-conduc-tor IEEE-Xplore}

\newcolumntype{Y}{>{\RaggedRight\arraybackslash}X}
\newcolumntype{P}[1]{>{\RaggedRight\arraybackslash}p{#1}}

\newcommand{\status}[1]{\textbf{#1}}

\begin{document}
	
	\title{LAB-DST SOFIA O26}
	
	\author{
		Marco Heimdal Castro Magaña, Alejandro Amador Franco, Karen Alejandra Guevara Vasquez, María Alicia Vela Contreras, Juan Pablo Contreras Quiñones,
		Jorge A. Lizarraga and Luis E. Gonzales Jimenez 
	}
	
	\markboth{Laboratory for the Development of Technological Solutions}
	{Technical Datasheet of SOFIA Demonstrator Module}
	
	\maketitle
	
	\begin{abstract}
		The abstract is a single paragraph of 150--250 words, written last,
		once results already exist. It must be an objective synthesis of the
		whole paper, without stating conclusions that are not supported in the
		body of the text. Even though it has no visible subheadings, draft it
		following this internal structure: \\
		(1) \textbf{Context}: what question or need motivates the work? \\
		(2) \textbf{Methods}: what was done to address it? (approach,
		materials, main methods, in 1-2 lines). \\
		(3) \textbf{Results}: what was found? (the main findings, with
		concrete data if possible). \\
		(4) \textbf{Conclusion}: what does that finding imply? \\
		Avoid mathematical or chemical formulas inside the abstract.
	\end{abstract}
	
	\begin{IEEEkeywords}
		keyword 1, keyword 2, keyword 3 (3 to 6 terms specific to the
		topic of the paper, yet reasonably common within the discipline --
		they help other researchers find the work).
	\end{IEEEkeywords}
	
	\section{Introduction}
	
	% IEEE journal papers open the first paragraph with a drop-cap letter
	% via \IEEEPARstart. Keep this wrapper when you replace the guidance
	% text below with your actual first paragraph.
	\IEEEPARstart{S}{OFIA} is a modular robotic installation composed of 16
	mechatronic units arranged in a two-dimensional $4 \times 4$ array, with each
	unit providing two rotational degrees of freedom actuated by servomotors. The
	system aims to achieve coordinated control of the multiple actuators while
	enabling interaction through static and dynamic human gestures. For this
	purpose, SOFIA integrates embedded control, inter-module communication, and an
	external computer vision system for gesture recognition.
	
	The SOFIA system integrates three main capabilities: automatic configuration
	and communication of modules through a Modbus network, characterization of the
	mechanical time constant of the servomotors, and recognition of static and
	dynamic human gestures through computer vision. These capabilities require the
	integration of mechanical actuation, servo control, embedded processing, and
	communication subsystems. Since each of the 16 units uses two servomotors, the
	complete demonstrator requires the coordinated control of 32 actuators and the
	generation of multiple PWM signals. The embedded controller must coordinate
	the desired actuator positions, while the communication architecture should
	support the future integration of additional SOFIA modules without
	significantly increasing wiring and control complexity.
	
		\begin{table*}[!t]
		\centering
		\caption{Comparative matrix of related shape-changing interface systems}
		\label{tab:related_work_matrix}
		\scriptsize
		
		\setlength{\tabcolsep}{3pt}
		\renewcommand{\arraystretch}{1.1}
		
		\begin{tabularx}{\textwidth}
			{@{}P{2.2cm} P{3cm} P{3cm} P{4cm} Y@{}}
			
			\toprule
			Reference & Method / technology & Metric / focus & Main result & Gap with respect to SOFIA \\
			\midrule
			
			Relief \cite{leithinger2010relief} &
			Array of 120 motorized pins, scalable low-cost control &
			Actuation density, hardware/software scalability &
			Continuous surface capable of rendering and animating 3D shapes, with per-pin input sensing &
			Translational actuation only (1 DOF per pin); no independent rotation per unit \\
			
			inFORM \cite{follmer2013inform} &
			Pin-based shape display on a pre-existing platform &
			Dynamic physical affordances and constraints, object manipulation &
			Real-time direct-touch interaction and physical rendering of content &
			Same translational limitation as Relief; centralized architecture, not modular per unit \\
			
			TRANSFORM \cite{ishii2015transform} &
			Three large-scale pin displays, kinetic energy as input &
			Aesthetic transformation from static furniture to dynamic system &
			Over one thousand pins coordinated in real-time wave motion &
			Installation-scale system; no per-module orientable control or distributed communication between units \\
			
			PneUI \cite{yao2013pneui} &
			Pneumatically actuated multi-layer soft composite materials &
			Integrated input sensing and shape output in the material itself &
			Four functional applications (phicons, mobile, case, lamp) with controllable shape change &
			Soft/pneumatic actuation, not discrete rotational motion; difficult to scale to precise angular control \\
			
			LineFORM \cite{nakagaki2015lineform} &
			Serial chain of one-DOF servo actuators &
			Line/curve-based interaction affordances &
			Cord, mobile, and body-constraint interfaces through line-shaped actuation &
			Serial (chain) configuration, not a two-dimensional array; 1 DOF per unit, not 2 \\
			
			ChainFORM \cite{nakagaki2016chainform} &
			Identical modules with servo, touch sensing, visual output and embedded processing &
			Hardware modularity and reconfigurability &
			Modular, reconfigurable platform with rich per-module I/O &
			Modularity demonstrated, but still a linear chain, not a 2D array with 2 DOF per node \\
			
			\bottomrule
			
		\end{tabularx}
	\end{table*}
	
	Motorized arrays have been widely explored as a mechanism for generating
	dynamic physical surfaces. Relief introduced an actuated tabletop composed of
	120 individually addressable motorized pins and emphasized a scalable hardware
	and software architecture \cite{leithinger2010relief}. This concept was
	extended through inFORM, where shape actuation was employed to generate dynamic
	physical affordances and constraints and to manipulate physical objects
	\cite{follmer2013inform}. At a larger scale, TRANSFORM explored the
	transformation of static physical structures into dynamic systems driven by
	streams of data and energy \cite{ishii2015transform}. These approaches
	demonstrate the potential of coordinated arrays of actuators, although their
	architectures primarily focus on large pin-based or installation-scale physical
	surfaces.
	
	Alternative approaches have explored different actuation principles and
	modular configurations. PneUI demonstrated shape-changing interfaces based on
	pneumatically actuated soft composite materials, integrating sensing and active
	shape output within deformable structures \cite{yao2013pneui}. In contrast,
	LineFORM investigated electromechanical shape transformation through serial
	chains of one-degree-of-freedom servo actuators
	\cite{nakagaki2015lineform}. This concept evolved toward greater hardware
	modularity in ChainFORM, where identical modules incorporate motor actuation,
	sensing, visual output, and embedded processing
	\cite{nakagaki2016chainform}. These systems highlight modularity and
	distributed actuation as relevant design characteristics for reconfigurable
	physical interfaces.
	
	The main technological characteristics and differences identified among the
	reviewed shape-changing interfaces are summarized in
	Table~\ref{tab:related_work_matrix}.
	
	In contrast to the reviewed shape-changing interfaces, SOFIA combines a
	two-dimensional $4 \times 4$ arrangement with two rotational degrees of freedom
	per unit. Unlike pin-based displays, which predominantly constrain motion to a
	single translational axis, or line-based modular systems organized as serial
	chains, each SOFIA unit can independently modify its orientation while
	remaining part of the same physical surface. This configuration introduces
	design challenges related to coordinated actuation, modularity, scalability,
	control architecture, electrical integration, and communication. Based on the
	reviewed literature, these characteristics were identified as relevant criteria
	for evaluating the technological architecture of the demonstrator.
	Table~\ref{tab:related_work_matrix} summarizes the comparative analysis of the
	reviewed systems, while Table~\ref{tab:benchmarking} contrasts SOFIA with the
	reviewed solutions using the established comparison criteria.
	
	\begin{table}[H]
		\centering
		\caption{Benchmarking against related shape-changing systems}
		\label{tab:benchmarking}
		\scriptsize
		
		\setlength{\tabcolsep}{2pt}
		\renewcommand{\arraystretch}{1.05}
		
		\begin{tabularx}{\columnwidth}
			{@{}P{1.7cm} P{1.7cm} P{1cm} Y@{}}
			
			\toprule
			Solution & Configuration & DOF/ unit & Modularity / scalability\\
			\midrule
			
			inFORM \cite{follmer2013inform} &
			Pin array on fixed platform &
			1 (translational) &
			Low -- centralized control, not designed for modular expansion \\
			
			LineFORM / ChainFORM \cite{nakagaki2015lineform,nakagaki2016chainform} &
			Serial chain of modules &
			1 per module &
			Medium -- modular, but limited to a linear configuration \\
			
			PneUI \cite{yao2013pneui} &
			Pneumatic soft composite material &
			Continuous (non-discrete) &
			Low -- difficult to discretize into individually addressable nodes \\
			
			\textbf{SOFIA (proposed)} &
			2D array, 4$\times$4 modules &
			\textbf{2 (rotational)} &
			\textbf{High -- I2C bus + RS-485 designed for future module expansion} \\
			
			\bottomrule
			
		\end{tabularx}
	\end{table}
	
	From the functional decomposition of the demonstrator and its comparison with
	the literature, several technological opportunities were identified. One is
	the development of a scalable control architecture capable of independently
	coordinating all actuators while reducing the processing load on the main
	controller. Another is the implementation of an inter-module communication
	architecture that allows multiple SOFIA modules to operate as part of a larger
	system. Additional opportunities include the implementation of
	software-defined motion limits and smoothing strategies to prevent abrupt
	servo movements and reduce mechanical stress.
	
	Based on these opportunities, this work focuses on defining and justifying a
	preliminary electronic and control architecture for the SOFIA demonstrator.
	The main areas considered are actuator control, PWM generation, embedded
	processing, inter-module communication, electrical compatibility, and
	scalability. These criteria provide the basis for the selection and
	experimental evaluation of the technologies used in the subsequent development
	stages of the SOFIA system.
	
	\section{Materials and Methods}
	\label{sec:methods}
	
	The SOFIA demonstrator consists of sixteen electromechanical units,
	each with two rotational degrees of freedom. The reference architecture
	requires thirty-two servo channels controlled by one ESP32-C3 and two
	PCA9685 PWM drivers \cite{sofia_architecture}.
	
	The objective of this document is to verify the compatibility of the
	selected components according to electrical, communication and
	mechanical requirements.
	The verification criteria are based on the hardware requirements and
	checklist defined in the technical session \cite{technical_session}.
	The document structure follows the LabDST documentation requirements for
	technical datasheets and component verification reports \cite{labdst_deliverables}.
	
	
	\begin{table}[H]
		\centering
		\caption{SOFIA Main Requirements}
		\scriptsize
		
		\setlength{\tabcolsep}{2pt}
		\renewcommand{\arraystretch}{1.05}
		
		\begin{tabularx}{\columnwidth}
			{@{}P{2.1cm} P{2.1cm} Y@{}}
			
			\toprule
			Requirement & Specification & Verification\\
			\midrule
			
			Power &
			5 V servo supply &
			Compatible with MG90S range \\
			
			Logic voltage &
			3.3 V MCU logic &
			Required for ESP32-C3 GPIO \\
			
			Communication &
			WiFi/Bluetooth + I2C &
			Supported by ESP32-C3 \\
			
			PWM control &
			32 servo channels, 50 Hz &
			Two PCA9685 drivers required \\
			
			Mechanical &
			16 modules, 2 DOF each &
			32 servos required \\
			
			\bottomrule
			
		\end{tabularx}
	\end{table}
	
	
	\subsection{ESP32-C3 Technical Analysis}
	
	
	\begin{table}[H]
		\centering
		\caption{ESP32-C3 Main Specifications}
		\scriptsize
		
		\setlength{\tabcolsep}{2pt}
		\renewcommand{\arraystretch}{1.05}
		
		\begin{tabularx}{\columnwidth}
			{@{}P{2cm} P{2.2cm} Y@{}}
			
			\toprule
			Parameter & Value & Verification\\
			\midrule
			
			Supply voltage &
			3.0--3.6 V &
			Requires regulated 3.3 V supply \\
			
			WiFi &
			802.11 b/g/n 2.4 GHz &
			Complies \\
			
			Bluetooth &
			BLE 5 &
			Complies \\
			
			Logic level &
			3.3 V GPIO &
			5 V cannot be applied directly \\
			
			Boot pins &
			GPIO2, GPIO8, GPIO9 &
			PCB design verification required \\
			
			Temperature &
			Up to 105 $^\circ$C &
			Suitable for laboratory operation \\
			
			\bottomrule
			
		\end{tabularx}
		
	\end{table}
	
	
	Source: ESP32-C3 Datasheet, Espressif Systems. \cite{esp32c3}
	
	
	\subsection{PCA9685 Technical Analysis}
	
	
	\begin{table}[H]
		\centering
		\caption{PCA9685 Main Specifications}
		\scriptsize
		
		\setlength{\tabcolsep}{2pt}
		\renewcommand{\arraystretch}{1.05}
		
		\begin{tabularx}{\columnwidth}
			{@{}P{2cm} P{2.2cm} Y@{}}
			
			\toprule
			Parameter & Value & Verification\\
			\midrule
			
			Channels &
			16 PWM channels &
			Two devices provide 32 outputs \\
			
			Resolution &
			12 bits &
			Complies \\
			
			Interface &
			I2C Fast Mode Plus &
			Compatible with ESP32-C3 \\
			
			Voltage &
			2.3--5.5 V &
			Use 3.3 V logic \\
			
			Frequency &
			24--1526 Hz &
			Configure to 50 Hz for servos \\
			
			Addressing &
			Up to 62 devices &
			Use different addresses (0x40/0x41) \\
			
			\bottomrule
			
		\end{tabularx}
		
	\end{table}
	
	
	Source: PCA9685 Datasheet, NXP. \cite{pca9685}
	
	
	\subsection{MG90S Servo Technical Analysis}
	
	
	\begin{table}[H]
		\centering
		\caption{MG90S Main Specifications}
		\scriptsize
		
		\setlength{\tabcolsep}{2pt}
		\renewcommand{\arraystretch}{1.05}
		
		\begin{tabularx}{\columnwidth}
			{@{}P{2cm} P{2.2cm} Y@{}}
			
			\toprule
			Parameter & Value & Verification\\
			\midrule
			
			Voltage & 4.8--6.0 V & Compatible with 5 V rail \\
			
			Torque & 1.8--2.2 kgf cm &  validation required \\
			
			Speed & 0.08--0.10 s/60$^\circ$ & Reference value \\
			
			Mass & 13.4 g & Must be considered in CAD \\
			
			PWM frequency & 50 Hz & Compatible with PCA9685 \\
			
			Stall current & Not reported & Measurement required \\
			
			\bottomrule
			
		\end{tabularx}
		
	\end{table}
	
	
	Source: TowerPro MG90S documentation. \cite{mg90s}
	
	\subsection{Checklist Verification}
	
	
	\begin{table}[H]
		\centering
		\caption{Verification Summary}
		\scriptsize
		
		\setlength{\tabcolsep}{2pt}
		\renewcommand{\arraystretch}{1.05}
		
		\begin{tabularx}{\columnwidth}
			{@{}P{2cm} P{1.7cm} Y@{}}
			
			\toprule
			Component &
			Status &
			Comment\\
			\midrule
			
			ESP32-C3 &
			\status{Complies} &
			Requires 3.3 V regulation and boot verification \\
			
			PCA9685 &
			\status{Complies} &
			Requires unique I2C addresses and 50 Hz configuration \\
			
			MG90S &
			\status{Partial} &
			Torque and current require validation \\
			
			\bottomrule
			
		\end{tabularx}
		
	\end{table}
	
	\subsection{Bill of Materials}
	
	\begin{table}[H]
		\centering
		\caption{Bill of materials summary}
		\label{tab:bom}
		\begin{tabularx}{\columnwidth}
			{@{}P{1.cm} P{3.5cm} P{0.5cm} Y@{} }
			\hline
			\textbf{Item} & \textbf{Component} & \textbf{Qty} & \textbf{Manufacturer} \\
			\hline
			1 & Sofia003-P001-03 & 16 & 3D Printed \\
			2 & MG90S servo & 32 & Tower Pro \\
			3 & Sofia003-E002-01 & 16 & 3D Printed \\
			3.1/4.4 & Pro3 & 32 & Tower Pro\\
			3.2 & Sofia003-P004-02 & 16 & 3D Printed \\
			3.3 & Sofia003-P004-03 & 16 & 3D Printed \\
			3.4/3.5/
			\par 4.3/7/8 & M2 screws (various) & 160 & Walfront \\
			3.6 & Axis 1 9mm x 9mm & 16 & Mae \\
			3.7 & Sofia003-P002-02 & 16 & 3D Printed \\
			4 & Sofia003-E003-01 & 16 & 3D Printed \\
			4.1 & Sofia003-P005-01 & 16 & 3D Printed \\
			4.2 & Sofia003-P005-02 & 16 & 3D Printed \\
			5 & Sofia003-P008-01 & 16 & 3D Printed \\
			6 & Sofia003-P009-01 & 16 & 3D Printed \\
			9 & Inser M3x5x4 & 16 & ZWMSSLL \\
			10 &Sofia003-P000-01 & 16 & 3D Printed \\
			- & ESP32-C3 SuperMini & 1 & Espressif \\
			- & PCA9685 Servo Driver & 2 & NXP \\
			- & MAX13487E TTL485-V2.0 & 16 & Analog Devices Inc.\\
			- & LED RGB & 1 & Mexbit \\
			- & JST-PH Male & 1 & JTAREA \\
			- & JST-PH Female & 1 & JTAREA \\
			
			\hline
		\end{tabularx}
	\end{table}
	
	\subsection{Justification of Technological Selection}
	
	\begin{table}[H]
		\centering
		\caption{Comparative Matrix of Alternatives}
		\label{tab:comparative_matrix}
		\scriptsize
		
		\setlength{\tabcolsep}{2pt}
		\renewcommand{\arraystretch}{1.05}
		
		\begin{tabularx}{\columnwidth}
			{@{}P{2cm} P{1.7cm} P{2cm} Y@{} }
			
			\toprule
			Function & Alternatives considered & Selected & Selection criterion \\
			\midrule
			
			Actuation & Servo / Stepper / DC & MG90S & Integrated angular positioning \\
			
			PWM generation & MCU PWM / PCA9685 / Additional MCU & PCA9685 & 16 PWM channels through I2C \\
			
			Processing & Arduino / STM32 / ESP32-C3 SuperMini & ESP32-C3 SuperMini & Wireless connectivity + required interfaces \\
			
			Inter-module communication & UART / CAN/ RS-485 & RS-485 & Differential multi-node communication \\
			
			
			\bottomrule
			
		\end{tabularx}
		
	\end{table}
	
	As summarized in Table~\ref{tab:comparative_matrix}, the technological
	architecture was selected according to the functional and scalability
	requirements of the SOFIA demonstrator. MG90S servomotors were selected
	to provide controlled angular positioning for each degree of freedom,
	while two PCA9685 drivers provide the 32 PWM channels required by the
	system through a shared I2C bus. The ESP32-C3 SuperMini was selected as
	the main controller due to its integrated wireless connectivity and
	support for the required communication interfaces. Finally, RS-485 was
	selected for inter-module communication because its differential and
	multi-node architecture supports the future connection of multiple
	SOFIA modules.
	
	\section{Design and Implementation of the Technological Solution}
	\label{sec:design}
	
	Figure~\ref{fig1} shows the simulated response of the sixteen-module SOFIA structure while tracking the moving point. The sequence illustrates each module reorientation at different instants, where each module adjusts its two rotational degrees of freedom to maintain alignment with  reference position.
	
	\begin{figure}[H]
		\centering
		\subfloat[Step 1]{\includegraphics[width=1.1in]{images/Sim_a}}
		\hfil
		\subfloat[Step 2]{\includegraphics[width=1.1in]{images/Sim_b}}
		\hfil
		\subfloat[Step 3]{\includegraphics[width=1.1in]{images/Sim_c}}
		\hfil
		\subfloat[Step 4]{\includegraphics[width=1.1in]{images/Sim_d}}
		\hfil
		\subfloat[Step 5]{\includegraphics[width=1.1in]{images/Sim_e}}
		\hfil
		\subfloat[Step 6]{\includegraphics[width=1.1in]{images/Sim_f}}
		\caption{Simulation of the SOFIA module tracking a moving reference point using inverse kinematics.}
		\label{fig1}
	\end{figure}
	
	Algorithm~\ref{alg:sofia_tracking} summarizes the inverse kinematics and smoothing strategy used by each SOFIA module to orient its platform toward the moving target.
	
	\begin{algorithm}[H]
		\caption{SOFIA module target tracking control}
		\label{alg:sofia_tracking}
		
		\begin{algorithmic}[1]
			
			\Require Target position $p$, module position $r_0$, joint limits $q_{1min},q_{1max},q_{2min},q_{2max}$, smoothing gain $K$
			\Ensure Servo angles $q_1,q_2$
			
			\State $v \gets p-r_0$
			\State $\hat{v} \gets (v/\|v\|)$
			
			\State $q_1^* \gets atan2(-\hat{v}_y,\hat{v}_z)$
			\State $q_2^* \gets asin(\hat{v}_x)$
			
			\State $q_1^* \gets Clamp(q_1^*,q_{1min},q_{1max})$
			\State $q_2^* \gets Clamp(q_2^*,q_{2min},q_{2max})$
			
			\State $q_1 \gets q_{1old}+K(q_1^*-q_{1old})$
			\State $q_2 \gets q_{2old}+K(q_2^*-q_{2old})$
			
			\State Apply $q_1,q_2$ to servo motors
			
			\State \Return $q_1,q_2$
			
		\end{algorithmic}
	\end{algorithm}
	
	
	
	
	\balance
	
	\bibliographystyle{IEEEtran}
	\bibliography{bibliography}
	
	
%	\begin{table*}[!t]
%		\centering
%		\caption{Bill of materials for one SOFIA module}
%		\label{tab:bom_f}
%		\begin{tabular}{p{3cm}cp{6cm}p{3cm}p{3cm}}
%			\hline
%			\textbf{Component} & \textbf{Qty} & \textbf{Technical description} &
%			\textbf{Manufacturer} & \textbf{Country} \\
%			\hline
%%			Continuous-rotation servo, rated torque $\tau$ at $V$~V &
%			Manufacturer Y & City, Country \\
%			Module enclosure & 16 &
%			3D-printed enclosure, PLA, in-house design &
%			In-house & City, Country \\
%			\hline
%		\end{tabular}
%	\end{table*}
	
\end{document}